\section{Introduction}
\label{sec:introduction}

Post-training unlearning aims to suppress targeted knowledge or training-data
influence while preserving behavior elsewhere
\citep{yao2024machineunlearning,li2024wmdp}. Comparisons with the original
model measure suppression and preservation, but do not show whether an edited
model resembles one trained without the target learning signal. We ask whether
post-training erasure can reproduce the behavior of a matched model trained
without loss on concept-associated chunks while preserving knowledge on neighboring
concepts and general-science behavior.

Using LMEnt's entity-annotated corpus, we construct matched references for
Ancient Rome, Baseball, and artificial intelligence
\citep{gottesman2025lment}. Alongside a full-corpus model, we train one
concept-excluded twin per concept with the same initialization and training
procedure, masking loss on linked chunks. We apply EMBER, RMU, and an
SNMF-based MLP erasure to copies of the full model and also evaluate RMU+EMBER
and SNMF+EMBER ensembles
\citep{suslik2026ember,li2024wmdp,shafran2026snmf}. Each edit is compared with
both the full model and its matched twin.

We evaluate target and neighboring questions with a shared SciQ
general-science set. Our primary metrics are a harmonic efficacy--preservation
score, normalized absolute question-level NLL distance from the twin, and
normalized teacher-forced full-vocabulary KL from the twin. Signed mean NLL
proximity and the frequency of question-level improvement provide
complementary diagnostics.

No method is uniformly best. RMU has the lowest absolute NLL distance for
Rome; SNMF has the lowest distance for Baseball and AI and the lowest
normalized KL throughout, but produces little target erasure. EMBER has the
largest target-specific NLL changes among standalone methods and the highest
efficacy--preservation score for Baseball; RMU leads for Rome. The selected
RMU+EMBER condition leads for AI, although its RMU component fails the
activity check. The highest-scoring method is closest under absolute NLL only
once and never has the lowest normalized KL. Target suppression with preserved
neighboring and SciQ accuracy therefore does not imply similarity to a matched
concept-excluded twin in these experiments.

Our contributions are threefold: we construct matched concept-excluded
references for evaluating post-training erasure; compare five erasure
conditions using separate efficacy, question-level similarity, and
distributional-similarity criteria; and show that strong target suppression
does not generally imply behavioral similarity to training without the
target-associated data.
